%========================================
% MatinBook — Stage 13: Index Test
% Version: v1.1
% Purpose:
%   Validate the index system of MatinBook v1.1:
%     1. Main index entries (\index{term})
%     2. Sub-entries (\index{term!sub})
%     3. Sub-sub-entries (\index{term!sub!subsub})
%     4. Bold entries (\idxbold)
%     5. See references (\idxsee)
%     6. See also references (\idxseealso)
%     7. xindy with Persian sorting (persian-variant2)
%     8. \printindex in the back matter
%
% Compilation:
%   xelatex -shell-escape stage13-index.tex
%   xelatex -shell-escape stage13-index.tex
%   xelatex -shell-escape stage13-index.tex
%   (xindy is called automatically via -shell-escape;
%    3 passes recommended for a stable index)
%
% Expected outcome:
%   A professionally typeset PDF demonstrating that the
%   index is generated correctly with Persian sorting,
%   sub-entries are properly nested, and see/see-also
%   references resolve to the correct entries.
%========================================

%------------------------------------------------------------------------
% Search Path (for tests directory)
%------------------------------------------------------------------------
\makeatletter
\def\input@path{{../}}
\makeatother

%------------------------------------------------------------------------
% Document Class
%------------------------------------------------------------------------
\documentclass{matinbook}

%------------------------------------------------------------------------
% Book Metadata
%------------------------------------------------------------------------
\title{آزمون سیستم نمایه}
\author{تیم توسعه MatinBook}
\date{مهر ۱۴۰۵}

\booktitle{آزمون سیستم نمایه}

\renewcommand{\repository}{github.com/matinmahmoudi-cs/matinbook}

%========================================================================
% DOCUMENT
%========================================================================
\begin{document}

%------------------------------------------------------------------------
% FRONT COVER
%------------------------------------------------------------------------
\makecover
    {آزمون سیستم نمایه}
    {ارزیابی مدخل، زیرمدخل و ارجاع در MatinBook}
    {تیم توسعه MatinBook}
    {مهر ۱۴۰۵}

%------------------------------------------------------------------------
% FRONT MATTER (symmetric margins)
%------------------------------------------------------------------------
\frontmatter
\frontmattergeometry

% Title page
\maketitle

% Copyright / Colophon
\thispagestyle{empty}
\vfill
\begin{center}
    {\large\textbf{شناسنامه آزمون}}

    \vspace{0.8cm}

    \begin{tabular}{rl}
        عنوان: & آزمون سیستم نمایه \\
        نسخه: & \matinbookversion \\
        تاریخ: & \matinbookdate \\
        موتور: & \lr{XeLaTeX} \\
        مجوز: & \lr{MIT} \\
    \end{tabular}

    \vspace{0.8cm}

    {\small این سند بخشی از مجموعه آزمون‌های یکپارچگی \lr{MatinBook} است.}
\end{center}
\clearpage

% Preface
\chapter*{پیشگفتار}
\addcontentsline{toc}{chapter}{پیشگفتار}

این سند، سیزدهمین آزمون از مجموعه‌ی آزمون‌های یکپارچگی
\lr{MatinBook} نسخه‌ی \lr{\matinbookversion} است. هدف آن،
اعتبارسنجی کامل سیستم نمایه‌ی کلاس در هشت حوزه‌ی متمایز
است: مدخل اصلی، زیرمدخل، زیرزیرمدخل، مدخل با فونت
سیاه، ارجاع «نگاه کنید به»، ارجاع «همچنین نگاه کنید
به»، مرتب‌سازی فارسی با \lr{xindy}، و چاپ نمایه با
\lr{\textbackslash printindex}.

در کتاب‌های فنی و علمی، نمایه یکی از ابزارهای مهم
ناوبری است که به خواننده کمک می‌کند مفاهیم مورد نظر
را سریع پیدا کند. یک نمایه‌ی حرفه‌ای باید با
مرتب‌سازی دقیق، زیرمدخل‌های سلسله‌مراتبی و ارجاع‌های
متقابل، تجربه‌ی خواندن را بهبود ببخشد. \lr{MatinBook}
با استفاده از \lr{xindy} و ماژول \lr{persian-variant2}،
این اهداف را محقق می‌سازد.

هر بخش از این آزمون، با مثال‌های عملی و ارجاعات متقابل
همراه است تا خواننده بتواند درستی رفتار هر زیرسیستم را
در بافت واقعی یک کتاب ارزیابی کند.

\clearpage

% Table of Contents
\tableofcontents
\clearpage

% List of Figures
\listoffigures
\clearpage

% List of Tables
\listoftables
\clearpage

%------------------------------------------------------------------------
% MAIN MATTER (wide margin for notes)
%------------------------------------------------------------------------
\mainmatter
\mainmattergeometry

%========================================================================
% CHAPTER 1 — PROGRAMMING CONCEPTS
%========================================================================
\chapter{مفاهیم برنامه‌نویسی}
\label{chap:programming}

\section{زبان‌های برنامه‌نویسی}
\label{sec:programming-languages}

\subsection{پایتون}
\label{sec:programming-python}

پایتون\index{پایتون} یک زبان برنامه‌نویسی سطح بالا،
تفسیری و همه‌منظوره است. این زبان توسط خیدو فان
روسوم\index{فان روسوم، خیدو} در سال ۱۹۹۱ ساخته شد.

\mnote{پایتون در سال ۲۰۲۴ محبوب‌ترین زبان برنامه‌نویسی جهان بوده است.}

ویژگی‌های اصلی پایتون\index{پایتون!ویژگی‌ها} عبارت‌اند از:

\begin{itemize}
    \item \textbf{خوانایی بالا:} نحو ساده و شبیه به
    زبان طبیعی\index{پایتون!نحو}
    \item \textbf{کتابخانه‌ی غنی:} کتابخانه‌ی استاندارد
    گسترده\index{پایتون!کتابخانه استاندارد}
    \item \textbf{چندسکویی:} اجرا روی سیستم‌عامل‌های
    مختلف\index{چندسکویی}
\end{itemize}

\subsection{زبان C++}
\label{sec:programming-cpp}

\lr{C++}\index{C++} یک زبان برنامه‌نویسی سطح میانی
با قابلیت‌های شیءگرایی\index{شیءگرایی} است. این زبان
توسط بیارنه استراستروپ\index{استراستروپ، بیارنه} در
سال ۱۹۸۵ توسعه یافت.

ویژگی‌های \lr{C++}\index{C++!ویژگی‌ها}:

\begin{itemize}
    \item \textbf{سرعت بالا:} کامپایل مستقیم به کد
    ماشین\index{C++!سرعت}
    \item \textbf{مدیریت حافظه:} کنترل دستی حافظه
    \index{مدیریت حافظه}
    \item \textbf{برنامه‌نویسی جنریک:} قالب‌ها
    \index{برنامه‌نویسی جنریک}
\end{itemize}

\section{الگوریتم‌ها}
\label{sec:programming-algorithms}

الگوریتم\index{الگوریتم} مجموعه‌ای از دستورالعمل‌های
گام‌به‌گام برای حل یک مسئله است.

\subsection{انواع الگوریتم‌ها}
\label{sec:programming-algorithms-types}

\begin{itemize}
    \item \textbf{جستجو:} جستجوی خطی\index{الگوریتم!جستجوی خطی}
    و جستجوی دودویی\index{الگوریتم!جستجوی دودویی} از
    الگوریتم‌های پایه‌ی جستجو هستند.

    \item \textbf{مرتب‌سازی:} مرتب‌سازی سریع\index{الگوریتم!مرتب‌سازی سریع}
    و مرتب‌سازی ادغامی\index{الگوریتم!مرتب‌سازی ادغامی}
    از پرکاربردترین الگوریتم‌های مرتب‌سازی هستند.

    \item \textbf{گراف:} الگوریتم دایکسترا\index{الگوریتم!دایکسترا}
    برای یافتن کوتاه‌ترین مسیر\index{کوتاه‌ترین مسیر}
    در گراف‌های وزن‌دار\index{گراف!وزن‌دار} استفاده
    می‌شود.
\end{itemize}

\section{ساختمان داده‌ها}
\label{sec:programming-datastructures}

ساختمان داده\index{ساختمان داده} روشی برای
سازمان‌دهی و ذخیره‌سازی داده‌ها در کامپیوتر است.

انواع ساختمان داده\index{ساختمان داده!انواع}:

\begin{itemize}
    \item \textbf{آرایه}\index{آرایه}: مجموعه‌ای از
    عناصر با اندیس\index{اندیس} عددی
    \item \textbf{لیست پیوندی}\index{لیست پیوندی}:
    مجموعه‌ای از گره‌ها\index{گره} که هر کدام به عنصر
    بعدی اشاره می‌کنند
    \item \textbf{پشته}\index{پشته}: ساختار
    \lr{LIFO}\index{LIFO} (آخرین ورودی، اولین خروجی)
    \item \textbf{صف}\index{صف}: ساختار
    \lr{FIFO}\index{FIFO} (اولین ورودی، اولین خروجی)
    \item \textbf{درخت}\index{درخت}: ساختار
    سلسله‌مراتبی\index{ساختار سلسله‌مراتبی} با
    ریشه\index{ریشه} و برگ‌ها\index{برگ}
    \item \textbf{گراف}\index{گراف}: مجموعه‌ای از
    رئوس\index{رأس} و یال‌ها\index{یال}
    \item \textbf{جدول درهم‌سازی}\index{جدول درهم‌سازی}:
    ساختار کلید-مقدار\index{کلید-مقدار} با دسترسی
    \lr{$O(1)$}
\end{itemize}

%========================================================================
% CHAPTER 2 — MATHEMATICAL CONCEPTS
%========================================================================
\chapter{مفاهیم ریاضی}
\label{chap:mathematics}

\section{نظریه‌ی اعداد}
\label{sec:math-numbertheory}

نظریه‌ی اعداد\index{نظریه اعداد} شاخه‌ای از
ریاضیات\index{ریاضیات} است که به مطالعه‌ی اعداد
صحیح\index{اعداد صحیح} می‌پردازد.

\subsection{اعداد اول}
\label{sec:math-primes}

عدد اول\index{عدد اول} عددی طبیعی بزرگ‌تر از ۱ است
که فقط بر ۱ و خودش بخش‌پذیر باشد.

قضیه‌ی اساسی حساب\index{قضیه!اساسی حساب} بیان
می‌کند که هر عدد طبیعی بزرگ‌تر از ۱ به‌صورت
حاصل‌ضرب اعداد اول\index{اعداد اول!تجزیه} قابل
نمایش است.

\subsection{بزرگ‌ترین مقسوم‌علیه مشترک}
\label{sec:math-gcd}

بزرگ‌ترین مقسوم‌علیه مشترک\index{GCD} یا \lr{GCD}
دو عدد، بزرگ‌ترین عددی است که هر دو را می‌شمارد.

الگوریتم اقلیدس\index{اقلیدس!الگوریتم} روشی کارآمد
برای محاسبه‌ی \lr{GCD} است.

\section{جبر}
\label{sec:math-algebra}

جبر\index{جبر} شاخه‌ای از ریاضیات است که به مطالعه‌ی
ساختارهای جبری\index{ساختار جبری} می‌پردازد.

\subsection{گروه}
\label{sec:math-group}

گروه\index{گروه} یک ساختار جبری شامل یک مجموعه و یک
عمل دوتایی\index{عمل دوتایی} با چهار ویژگی بسته
بودن\index{گروه!بسته بودن}، شرکت‌پذیری\index{گروه!شرکت‌پذیری}،
وجود عضو خنثی\index{گروه!عضو خنثی} و وجود عضو
وارون\index{گروه!عضو وارون} است.

\subsection{ماتریس}
\label{sec:math-matrix}

ماتریس\index{ماتریس} آرایه‌ای مستطیلی از اعداد است.
ماتریس‌ها\index{ماتریس!انواع} در جبر خطی\index{جبر خطی}،
گرافیک کامپیوتری\index{گرافیک کامپیوتری} و یادگیری
ماشین\index{یادگیری ماشین} کاربرد دارند.

\section{آنالیز ریاضی}
\label{sec:math-analysis}

آنالیز ریاضی\index{آنالیز ریاضی} به مطالعه‌ی
حد\index{حد}، پیوستگی\index{پیوستگی}،
مشتق\index{مشتق} و انتگرال\index{انتگرال} می‌پردازد.

قضیه‌ی مقدار میانی\index{قضیه!مقدار میانی} و
قضیه‌ی اساسی حسابان\index{قضیه!اساسی حسابان} از
قضایای مهم آنالیز هستند.

%========================================================================
% CHAPTER 3 — USING INDEX COMMANDS
%========================================================================
\chapter{راهنمای استفاده از نمایه}
\label{chap:index-commands}

\section{دستورات پایه}
\label{sec:index-basic}

برای ایجاد مدخل در نمایه از دستور
\lr{\textbackslash index} استفاده می‌شود:

\begin{itemize}
    \item \lr{\textbackslash index\{واژه\}} —
    مدخل اصلی\index{مدخل اصلی}
    \item \lr{\textbackslash index\{واژه!زیرواژه\}} —
    زیرمدخل\index{زیرمدخل}
    \item \lr{\textbackslash index\{واژه!زیرواژه!زیرزیرواژه\}} —
    زیرزیرمدخل\index{زیرزیرمدخل}
\end{itemize}

\section{دستورات پیشرفته}
\label{sec:index-advanced}

\lr{MatinBook} دستورات کمکی زیر را برای نمایه
فراهم می‌کند:

\begin{itemize}
    \item \lr{\textbackslash idxbold}: مدخل با
    فونت سیاه\idxbold{سیاه}
    \item \lr{\textbackslash idxsee}: ارجاع «نگاه
    کنید به»\index{ارجاع!نگاه کنید به}
    \item \lr{\textbackslash idxseealso}: ارجاع
    «همچنین نگاه کنید به»\index{ارجاع!همچنین نگاه کنید به}
    \item \lr{\textbackslash idxsub}: ایجاد زیرمدخل
    سریع
    \item \lr{\textbackslash idxsubsub}: ایجاد
    زیرزیرمدخل سریع
\end{itemize}

\section{نمونه ارجاع‌های نمایه}
\label{sec:index-examples}

\begin{itemize}
    \item \idxsee{مرتب‌سازی سریع}{الگوریتم!مرتب‌سازی سریع}
    \item \idxseealso{درخت دودویی}{درخت}
    \item \idxbold{مفهوم مهم}
    \item \idxsub{برنامه‌نویسی}{پایتون}
\end{itemize}

\mnote{دستورات \lr{\textbackslash idxsee} و \lr{\textbackslash idxseealso} با \lr{xindy} سازگار هستند.}

%========================================================================
% CHAPTER 4 — SUMMARY
%========================================================================
\chapter{خلاصه‌ی نتایج}
\label{chap:summary}

\section{وضعیت سیستم نمایه}
\label{sec:summary-status}

جدول \cref{tab:index-status} وضعیت قابلیت‌های نمایه‌ی
\lr{MatinBook} را نشان می‌دهد.

\begin{matintable}{وضعیت قابلیت‌های نمایه}{tab:index-status}
\begin{tblr}{
    width = \textwidth,
    colspec = {l X[c] X[c]},
    hlines, vlines,
    colsep = 8pt,
    rowsep = 4pt,
    row{1} = {font=\bfseries},
}
قابلیت & توضیحات & وضعیت \\
مدخل اصلی & \lr{\textbackslash index\{واژه\}} & \textcolor{matingreen}{فعال} \\
زیرمدخل & \lr{\textbackslash index\{واژه!زیر\}} & \textcolor{matingreen}{فعال} \\
زیرزیرمدخل & \lr{\textbackslash index\{واژه!زیر!زیر\}} & \textcolor{matingreen}{فعال} \\
ارجاع \lr{see} & \lr{\textbackslash idxsee\{از\}\{به\}} & \textcolor{matingreen}{فعال} \\
ارجاع \lr{see also} & \lr{\textbackslash idxseealso\{از\}\{به\}} & \textcolor{matingreen}{فعال} \\
فونت سیاه & \lr{\textbackslash idxbold\{واژه\}} & \textcolor{matingreen}{فعال} \\
زیرمدخل سریع & \lr{\textbackslash idxsub\{اصلی\}\{زیر\}} & \textcolor{matingreen}{فعال} \\
چاپ نمایه & \lr{\textbackslash printindex} & \textcolor{matingreen}{فعال} \\
\end{tblr}
\end{matintable}

\section{معیارهای ارزیابی}
\label{sec:summary-criteria}

در این آزمون، سیستم نمایه‌ی \lr{MatinBook} بر اساس
پنج معیار ارزیابی شد:

\begin{enumerate}
    \item \textbf{مرتب‌سازی فارسی}: با \lr{xindy}.
    \item \textbf{زیرمدخل‌ها}: سلسله‌مراتب درست.
    \item \textbf{ارجاع‌ها}: \lr{see} و \lr{see also}.
    \item \textbf{فونت سیاه}: \lr{\textbackslash idxbold}.
    \item \textbf{چاپ نمایه}: در \lr{backmatter}.
\end{enumerate}

\section{نتیجه‌گیری}
\label{sec:summary-conclusion}

\textbf{تمام زیرسیستم‌های نمایه‌ی \lr{MatinBook}
نسخه‌ی \lr{\matinbookversion} با موفقیت بارگذاری و
آزمایش شدند. بیش از ۶۰ مدخل در دو فصل موضوعی با
زیرمدخل‌ها و ارجاع‌های متقابل، به‌درستی در نمایه
ثبت می‌شوند و با \lr{xindy} مرتب‌سازی فارسی
انجام می‌شود.}

%------------------------------------------------------------------------
% BACK MATTER (symmetric margins)
%------------------------------------------------------------------------
\backmatter
\frontmattergeometry

% Bibliography (empty placeholder)
\printbibliography[title={منابع و مراجع}]

% Index
\printindex

%------------------------------------------------------------------------
% BACK COVER
%------------------------------------------------------------------------
\authorbio{%
    تیم توسعه‌ی \lr{MatinBook}، گروهی از پژوهشگران و برنامه‌نویسان
    فارسی‌زبان است که با هدف ارائه‌ی یک راه‌حل حرفه‌ای برای
    حروف‌چینی کتاب‌های فنی فارسی فعالیت می‌کند. این پروژه حاصل
    سال‌ها تجربه در حوزه‌ی نشر دیجیتال و برنامه‌نویسی است.
}

\makebackcover{%
    این سند، سیزدهمین آزمون از مجموعه‌ی آزمون‌های یکپارچگی
    \lr{MatinBook} نسخه‌ی \lr{\matinbookversion} است.

    \vspace{0.5cm}

    \textbf{زیرسیستم‌های آزموده‌شده:}
    \begin{itemize}
        \item مدخل اصلی
        \item زیرمدخل و زیرزیرمدخل
        \item مدخل با فونت سیاه
        \item ارجاع \lr{see}
        \item ارجاع \lr{see also}
        \item مرتب‌سازی فارسی با \lr{xindy}
        \item چاپ نمایه در \lr{backmatter}
    \end{itemize}
}

\end{document}